问题二的困难在于，日前必须锁定各时段的普通购电额度，而负荷和光伏偏差会在日内逐步揭示。若只按均值预测购电，正净负荷偏差可能迫使系统以 $5$ 倍价格紧急购电；若机械地扩大计划额度，又会因未取用额度仍按计划价格结算而增加成本。因此，预测的评价不能只取负荷或光伏的 MAE，而应以实际计划费和紧急购电费之和衡量。

本文把一天分为“日前计划”和“日内执行”两个信息阶段。每日 $0{:}00$，仅由此前已经结束的日期构造负荷、光伏预测和残差池，并据此锁定唯一的普通购电额度 $q_{d,t}$；在每个 $10$ 分钟时段结束后，实际负荷、光伏和 SOC 才进入信息集，控制器仅重算下一步储能动作。由此，历史残差可用于描述未来偏差的形态，却不能作为当天已经确定的未来路径直接执行。

为同时保留负荷与光伏误差的时序结构并控制求解规模，本文以完整历史残差日为样本，经 K-medoids 选取代表轨迹；K 只作为场景压缩规模，不被解释为真实概率分布的精确离散化。日前层以残差分位数形成必要的净负荷储备，下限只作用于相应时段的 $q_{d,t}$，而不人为放大全日购电量。日内则依据已观测残差前缀对代表轨迹重新加权，配合滚动 MPC 决定充放电和紧急购电。

这一结构的关键是“政策一致性”：预测结构和风险分位的选择，也必须经历同样的日前锁定与逐时段 MPC 过程，不能把整日真实负荷、光伏交给事后全视域调度后再评分。这样，少买计划电和避免 $5$ 倍紧急购电由同一个货币目标自然权衡，结果可被解释为题设信息条件下的因果回测，而非预知未来的事后最优。
